Để cung cấp thông tin chi tiết về các tiến trình đang hoạt động trong hệ thống theo yêu cầu của system call \texttt{sysinfo}, thành viên C chịu trách nhiệm cài đặt bộ đếm lập lịch \texttt{nsched} và hàm thu thập danh sách tiến trình \texttt{proccollect()} trong \texttt{kernel/proc.c}.

\subsubsection{Cài đặt bộ đếm sự kiện lập lịch (\texttt{nsched})}
\begin{itemize}
    \item \textbf{Khai báo trường:} Trường \texttt{int nsched;} được bổ sung vào \texttt{struct proc} (\texttt{kernel/proc.h}) trong nhóm các trường dữ liệu được bảo vệ bởi khóa \texttt{p->lock}.
    \item \textbf{Khởi tạo:} Trong hàm \texttt{allocproc()} (\texttt{kernel/proc.c}), ngay sau khi slot tiến trình được gán trạng thái \texttt{USED}, biến được khởi tạo về $0$:
    \begin{lstlisting}[language=C]
    found:
      p->pid = allocpid();
      p->state = USED;
      p->nsched = 0;
    \end{lstlisting}
    \item \textbf{Tăng giá trị:} Trong hàm \texttt{scheduler()} (\texttt{kernel/proc.c}), ngay trước lệnh \texttt{swtch(\&c->context, \&p->context)}, ta thực hiện tăng bộ đếm:
    \begin{lstlisting}[language=C]
      if(p->state == RUNNABLE) {
        p->state = RUNNING;
        c->proc = p;
        p->nsched++; // Tang so lan duoc scheduler chon chay
        swtch(&c->context, &p->context);
    \end{lstlisting}
    \item \textbf{Tính an toàn đồng bộ:} Thao tác tăng \texttt{p->nsched++} diễn ra khi CPU đang nắm giữ \texttt{p->lock}, do đó không bao giờ xảy ra race condition giữa các CPU khi cùng quan sát tiến trình.
\end{itemize}

\subsubsection{Cài đặt hàm thu thập \texttt{proccollect()}}
Hàm \texttt{int proccollect(uint64 addr)} có nhiệm vụ duyệt qua toàn bộ bảng \texttt{proc[0 \dots NPROC-1]}, thu thập các tiến trình có trạng thái khác \texttt{UNUSED}, và ghi dữ liệu ra không gian người dùng tại địa chỉ \texttt{addr}.

\begin{lstlisting}[language=C]
int
proccollect(uint64 addr)
{
  struct proc *p;
  struct sysinfo_proc sproc;
  int count = 0;
  struct proc *curproc = myproc();

  for(p = proc; p < &proc[NPROC]; p++){
    acquire(&wait_lock);
    acquire(&p->lock);

    if(p->state == UNUSED){
      release(&p->lock);
      release(&wait_lock);
      continue;
    }

    memset(&sproc, 0, sizeof(sproc));
    sproc.pid = p->pid;
    sproc.state = p->state;
    sproc.nsched = p->nsched;
    safestrcpy(sproc.name, p->name, sizeof(sproc.name));

    if(p->parent)
      sproc.ppid = p->parent->pid;
    else
      sproc.ppid = 0;

    int nfd = 0;
    for(int fd = 0; fd < NOFILE; fd++){
      if(p->ofile[fd]){
        sproc.fds[nfd] = fd;
        sproc.fidx[nfd] = fileidx(p->ofile[fd]);
        nfd++;
      }
    }
    sproc.nfd = nfd;

    release(&p->lock);
    release(&wait_lock);

    uint64 dst = addr + __builtin_offsetof(struct sysinfo, procs) 
                      + (uint64)count * sizeof(struct sysinfo_proc);
    if(copyout(curproc->pagetable, dst, (char *)&sproc, sizeof(sproc)) < 0)
      return -1;

    count++;
  }

  return count;
}
\end{lstlisting}

\subsubsection{Các giải pháp kỹ thuật cốt lõi trong \texttt{proccollect}}
\begin{enumerate}
    \item \textbf{Giải quyết bài toán tràn Kernel Stack (Stack Overflow):}
    Cấu trúc \texttt{struct sysinfo} theo thiết kế chứa mảng $100$ phần tử \texttt{sysinfo\_file} và $64$ phần tử \texttt{sysinfo\_proc}, với tổng kích thước lên đến $\sim 38.25\text{ KB}$. Trong khi đó, mỗi luồng kernel trong xv6 chỉ được cấp phát duy nhất một trang bộ nhớ ngăn xếp (\texttt{KSTACK = PGSIZE = 4KB}). Nếu khai báo một biến cục bộ \texttt{struct sysinfo info} trên stack kernel, hệ thống sẽ lập tức bị tràn trang bảo vệ (guard page) và panic ngay lập tức.
    
    \textit{Giải pháp:} Khai báo biến cục bộ duy nhất cho \textbf{từng phần tử} \texttt{struct sysinfo\_proc sproc} ($160$ bytes). Với mỗi tiến trình hợp lệ, dữ liệu được đóng gói vào \texttt{sproc} trên stack rồi lập tức sao chép ra không gian người dùng bằng \texttt{copyout()}. Địa chỉ đích được tính toán chính xác bằng macro trình biên dịch \texttt{\_\_builtin\_offsetof(struct sysinfo, procs) + count * sizeof(struct sysinfo\_proc)}.

    \item \textbf{Thứ tự khóa (Lock Ordering) chống Deadlock:}
    Trường \texttt{p->parent} được bảo vệ bởi spinlock \texttt{wait\_lock}. Trong xv6, ở các hàm quan trọng như \texttt{exit()} và \texttt{wait()}, thứ tự khóa luôn được quy định: \textbf{\texttt{wait\_lock} được lấy trước \texttt{p->lock}}. Nếu trong \texttt{proccollect} ta đảo ngược thứ tự (lấy \texttt{p->lock} trước rồi mới lấy \texttt{wait\_lock}), một tiến trình khác đang thực thi \texttt{exit()} (đã giữ \texttt{wait\_lock} và đang chờ lấy \texttt{p->lock}) sẽ gây ra tình trạng khóa chết vòng tròn (circular wait deadlock). Do đó, hàm tuân thủ nghiêm ngặt quy tắc: lấy \texttt{wait\_lock} trước, sau đó mới lấy \texttt{p->lock}.

    \item \textbf{Thả khóa trước khi thực hiện \texttt{copyout()}:}
    Hàm \texttt{copyout()} thực hiện truy xuất bảng phân trang và ghi dữ liệu vào không gian bộ nhớ của tiến trình người dùng, vốn là thao tác tiềm ẩn lỗi hoặc gây trễ. Trong xv6, việc giữ spinlock khi truy cập dữ liệu người dùng có thể dẫn đến vi phạm nghiêm trọng (làm tăng thời gian giữ khóa hoặc vi phạm các assertion về ngắt). Vì vậy, dữ liệu của tiến trình được chụp nhanh (snapshot) vào biến cục bộ \texttt{sproc}, toàn bộ khóa \texttt{p->lock} và \texttt{wait\_lock} được giải phóng ngay sau đó, rồi mới gọi \texttt{copyout()}.

    \item \textbf{Xử lý tiến trình gốc (\texttt{init}):}
    Tiến trình đầu tiên \texttt{initproc} không có cha (\texttt{p->parent == 0}). Hàm kiểm tra điều kiện con trỏ \texttt{p->parent} trước khi truy xuất, gán \texttt{ppid = 0} một cách an toàn để tránh dereference con trỏ NULL gây kernel trap.
\end{enumerate}
